\documentclass[12pt]{article}
\usepackage[utf8]{inputenc}
\usepackage[spanish]{babel}
\usepackage{amsmath}
\usepackage{amsfonts}
\usepackage{amssymb}
\usepackage{geometry}
\geometry{letterpaper, margin=2.5cm}

\begin{document}

\begin{center}
    \Large \textbf{Evaluación de Examen 1: Unidad I} \\
    \large Física para Ciencias de la Salud (005-1713) \\
    \vspace{0.5cm}
    \textbf{Estudiante:} Mariana Rojas \quad \textbf{C.I.:} 33.386.743
    \vspace{0.3cm} \\
    \large \textbf{Calificación Final: 9/10 puntos}
\end{center}

\vspace{0.5cm}

\section*{Análisis de Resultados}

\subsection*{1. Ángulo plano (Conversión a radianes)}
\textbf{Procedimiento y resultado:} Plantea correctamente la conversión multiplicando por el factor formal $\left(\frac{\pi}{180^\circ}\right)$. Simplifica la fracción de manera limpia $\left(\frac{75\pi}{180} = \frac{5\pi}{12}\right)$, obteniendo el valor exacto y acompañándolo además de su respectiva aproximación decimal ($0,42\pi$ rad). \\
\textbf{Veredicto:} \textbf{Correcto} (2/2 pts).

\subsection*{2. Sistemas de coordenadas (Longitud del hueso)}
\textbf{Procedimiento y resultado:} Excelente resolución analítica. Escribe la fórmula general de la distancia entre dos puntos y sustituye los valores paso a paso. Desarrolla las restas, eleva al cuadrado ($\sqrt{36+64} = \sqrt{100}$) y aplica correctamente la raíz cuadrada para hallar la longitud exacta de $10$ cm. \\
\textbf{Veredicto:} \textbf{Correcto} (2/2 pts).

\subsection*{3. Vectores (Fuerza resultante)}
\textbf{Procedimiento y resultado:} La estudiante evidencia un gran orden. Plantea formalmente que $\vec{F}_R = \vec{F}_1 + \vec{F}_2$. Agrupa las componentes $\hat{i}$ y $\hat{j}$ de forma aislada, respetando impecablemente la ley de los signos para ensamblar el vector final exacto: $(30\hat{i} + 60\hat{j})$ N. \\
\textbf{Veredicto:} \textbf{Correcto} (2/2 pts).

\subsection*{4. Potencias de diez (Notación científica y prefijo)}
\textbf{Procedimiento y resultado:} Expresa la medida de forma adecuada aplicando estrictamente las reglas de la notación científica ($1,25 \times 10^{-5}$ m). Sin embargo, omite la respuesta de la segunda parte requerida en el enunciado, la cual consistía en identificar e indicar el nombre del prefijo del Sistema Internacional (micrómetros). \\
\textbf{Veredicto:} \textbf{Incompleto / Parcialmente correcto} (1/2 pts).

\subsection*{5. Sistemas de unidades (Conversión de densidad)}
\textbf{Procedimiento y resultado:} Demuestra un dominio sobresaliente de los factores de conversión en cadena. Plantea la conversión del volumen de manera muy técnica, elevando la relación lineal al cubo $\left(\frac{100 \text{ cm}}{1 \text{ m}}\right)^3$. Desglosa las operaciones de las fracciones con orden hasta llegar al resultado numérico perfecto de $1030 \text{ kg/m}^3$. \\
\textbf{Veredicto:} \textbf{Correcto} (2/2 pts).

\vspace{0.5cm}
\section*{Comentarios Generales y Sugerencias}
¡Un examen excelente! La estudiante resalta por su orden, caligrafía clara y un nivel técnico superior en sus planteamientos matemáticos (como se observó al aplicar factores de conversión cúbicos de forma implícita).
\begin{itemize}
    \item La única sugerencia formativa es recordar leer minuciosamente cada enunciado para evitar perder puntos en instrucciones secundarias (como ocurrió al no nombrar el prefijo en la Pregunta 4).
\end{itemize}

\end{document}